\section{Threat Models for DNS Queries and Tor} \label{sec:threatmodel}

The threat model for anonymous \ac{dns} queries considers a passive adversary
with access to query information at the target recursive
resolver~\cite{RFC9230,kurihara2023muodns,DoHoT,ODNS}. The objective is to
anonymize the client's identity (IP address) so that it cannot be linked to the
issued query processed by the resolver. We are, however, not trying to hide the
fact that the client is communicating with a \ac{dns} resolver.

Oblivious approaches such as \ac{odoh} rely on a single-proxy model whose
security hinges entirely on a \emph{non-collusion assumption}: the proxy that
sees the client's IP must not cooperate with the resolver that sees the
decrypted query. If these entities collude, the client's privacy is immediately
compromised. This non-collusion assumption is hard to verify let alone to
enforce in practice, especially if both parties fall under the same operational
or legal domain.

In contrast, Tor considers a stronger threat model, where the adversary has a
\emph{partial} view and/or control of the relays and communication
links~\cite{dingledine2004} that make up the Tor network, consisting of over
9,000
relays\footnote{\url{https://metrics.torproject.org/networksize.html?start=2025-12-01}}.
In general, reliably linking user activity in Tor requires control of all relays
used for a circuit. 

For \ac{doh} queries over Tor, the entry relay can observe the client source IP
and encrypted traffic but not the \ac{dns} query itself. The middle relay only
sees encrypted traffic without access to either the client source IP, resolver
IP, or the query. The exit relay only sees encrypted traffic and the target
\ac{doh} resolver. The \ac{doh} resolver sees the decrypted \ac{dns} queries
and the exit source IP, but never the client source IP. Even if the Tor circuit
is shortened (e.g., by removing the middle relay), the design still ensures
that no single entity can observe both the client's identity and the content of
their queries simultaneously. Even if all relays are colluding in this setting,
they still do not have access to the query since it is encrypted when sent to
the \ac{doh} resolver. Thus, Tor's threat model remains stronger than that of
single oblivious proxies, which fail completely under collusion between the two
involved parties. 

